\documentclass[10pt]{article}
% Shared preamble and text for the data-request letters.
% Replace the bracketed fields before submitting any letter.

\usepackage[a4paper,top=2.4cm,bottom=2.4cm,left=2.8cm,right=2.4cm]{geometry}
\usepackage[utf8]{inputenc}
\usepackage[T1]{fontenc}
\usepackage{lmodern}
\usepackage{microtype}
\usepackage{setspace}
\usepackage{array}
\usepackage{booktabs}
\usepackage{enumitem}
\usepackage{tabularx}
\usepackage{xcolor}
\usepackage[hidelinks]{hyperref}

\setstretch{1.12}
\setlength{\parindent}{0pt}
\setlength{\parskip}{6pt}
\setlist[itemize]{leftmargin=1.5em,itemsep=2pt,topsep=2pt}
\setlist[enumerate]{leftmargin=1.7em,itemsep=2pt,topsep=2pt}
\newcolumntype{Y}{>{\raggedright\arraybackslash}X}
\newcolumntype{P}[1]{>{\raggedright\arraybackslash}p{#1}}

\newcommand{\NamaMahasiswa}{[NAMA LENGKAP MAHASISWA]}
\newcommand{\NIMMahasiswa}{[NIM]}
\newcommand{\NamaPembimbing}{[NAMA DAN GELAR DOSEN PEMBIMBING]}
\newcommand{\JudulSementara}{Diagnosis \textit{Borrowed Size} dan \textit{Agglomeration Shadow} antara Jakarta dan Pusat-pusat Sekunder di Kawasan Aglomerasi Jakarta}

\newcommand{\KopSurat}{%
  \begin{center}
    \fbox{%
      \begin{minipage}{0.88\textwidth}
        \centering
        \textbf{KOP SURAT RESMI FAKULTAS TEKNIK UNIVERSITAS GADJAH MADA}\\
        \small Ganti blok ini dengan kop resmi sebelum surat dikirim.
      \end{minipage}}
  \end{center}
  \vspace{4pt}
}

\newcommand{\IdentitasMahasiswa}{%
  \begin{tabularx}{\textwidth}{@{}P{0.30\textwidth}Y@{}}
    Nama & \NamaMahasiswa \\
    NIM & \NIMMahasiswa \\
    Program studi & Sarjana Perencanaan Wilayah dan Kota \\
    Departemen dan fakultas & Departemen Teknik Arsitektur dan Perencanaan, Fakultas Teknik, Universitas Gadjah Mada \\
    Dosen pembimbing & \NamaPembimbing
  \end{tabularx}
}

\newcommand{\PernyataanPenggunaanData}{%
Data yang dimohon akan digunakan untuk kepentingan pendidikan dan penelitian nonkomersial. Data tidak akan dijual, dialihkan, atau disebarluaskan dalam bentuk mentah. Hasil yang dipublikasikan akan berupa statistik agregat, peta turunan, atau keluaran lain yang diizinkan oleh pemilik data. Pengolahan dan penyimpanan data akan mengikuti ketentuan kerahasiaan, lisensi, pembatasan publikasi, dan perjanjian penggunaan data yang berlaku.
}

\newcommand{\PernyataanTanpaBiaya}{%
Apabila permintaan ini tidak dapat dipenuhi tanpa biaya, mohon kami diberi tahu terlebih dahulu. Kami tidak menyetujui pembuatan tabulasi atau pengadaan data berbayar tanpa persetujuan baru dari mahasiswa dan dosen pembimbing.
}

\newcommand{\AbstraksiPenelitian}{%
\textbf{Latar belakang.} Kegiatan ekonomi, perjalanan, layanan, dan perumahan di sekitar Jakarta berlangsung melampaui batas administrasi DKI Jakarta. Hubungan tersebut tidak otomatis menunjukkan bahwa pusat-pusat sekunder telah memperkuat fungsi lokalnya. Suatu wilayah dapat memperoleh akses ke pasar metropolitan, tetapi pada saat yang sama semakin bergantung pada Jakarta untuk pekerjaan, layanan, atau fungsi bernilai tinggi.

\textbf{Tujuan.} Penelitian ini mengkaji apakah integrasi dengan Jakarta berjalan bersama penguatan fungsi lokal, suatu pola yang konsisten dengan \textit{borrowed size}, atau bersama defisit fungsi dan ketergantungan yang asimetris, suatu pola yang konsisten dengan \textit{agglomeration shadow}. Penelitian tidak bertujuan mengusulkan perluasan wilayah DKI Jakarta, mengubah batas pemerintahan, atau membuktikan hubungan kausal hanya dari potret satu periode.

\textbf{Metode.} Kasus utama adalah sistem metropolitan Jabodetabek. Wilayah Statistik Metropolitan Indonesia digunakan sebagai sabuk pembanding. DKI Jakarta diperlakukan sebagai satu inti dengan batas administrasi tetap. Kecamatan menjadi unit utama jika resolusi data mendukung. Data pada tingkat kabupaten/kota, zona transportasi, raster, dan titik dipertahankan pada resolusi aslinya sebelum dilakukan agregasi yang dapat ditelusuri. Analisis membedakan data yang teramati, tabulasi agregat, hasil pemodelan, dan proksi.

\textbf{Cakupan wilayah dan waktu.} Baseline utama menggunakan tahun 2024. Data tahun lain digunakan sebagai konteks atau pemeriksaan perubahan, dengan tahun dan perbedaan unit ditampilkan secara terbuka. Wilayah mencakup DKI Jakarta, kabupaten/kota dalam kawasan utama Jabodetabek, pusat sekunder yang ditetapkan melalui prosedur yang terdokumentasi, serta sabuk pembanding yang diperlukan untuk menguji fungsi relatif.

\textbf{Jenis data dan variabel.} Penelitian memerlukan data penduduk, luas, kepadatan, integrasi dan perjalanan, fungsi layanan lokal, pekerjaan menurut lokasi tempat kerja dan sektor jika tersedia, aksesibilitas jaringan, serta karakteristik pusat sekunder. Podes digunakan untuk mengukur kehadiran dan akses layanan, bukan sebagai ukuran langsung produktivitas atau jumlah pekerjaan. Data komuter pada tingkat kabupaten/kota tidak akan dipaksa menjadi estimasi kecamatan. Data OD yang dimodelkan akan diberi label sintetis dan diuji terhadap kendala yang tersedia.

\textbf{Rencana hasil.} Hasil penelitian berupa profil integrasi dan fungsi lokal, tipologi kondisi relasional, peta gradien, uji sensitivitas bobot dan ambang, serta penjelajah skenario berbasis WebGL. Penjelajah tersebut menghitung ulang dan menampilkan hasil secara interaktif. Ia bukan sumber data waktu nyata, ramalan kebijakan, atau simulasi perubahan kewenangan pemerintahan.
}

\newcommand{\LampiranAbstraksi}{%
  \clearpage
  \section*{Lampiran 2. Abstraksi penggunaan data}
  \AbstraksiPenelitian
}

\newcommand{\TandaTanganDTAP}{%
  Hormat kami,\\[2.2cm]
  \parbox[t]{0.85\textwidth}{[NAMA PENANDATANGAN]}\\
  \parbox[t]{0.85\textwidth}{[JABATAN PENANDATANGAN SESUAI ALUR PERSURATAN DTAP]}
}

\newcommand{\TandaTanganBPS}{%
  Hormat kami,\\[2.2cm]
  \parbox[t]{0.85\textwidth}{[NAMA PENANDATANGAN]}\\
  \parbox[t]{0.85\textwidth}{[DEKAN FAKULTAS TEKNIK ATAU PEJABAT SETINGKAT YANG DITERIMA BPS]}
}

\usepackage{amssymb}
\usepackage{longtable}
\usepackage{pgffor}

\hypersetup{
  pdftitle={Skenario Akses Data Skripsi: Keluarga Metodologis dan Registri 64 Kombinasi},
  pdfauthor={\NamaMahasiswa}
}

\setlength{\parskip}{5pt}
\raggedbottom
\newcounter{combination}
\hypersetup{pageanchor=false}

\newcommand{\Door}[1]{\ifnum#1=1\textbf{1}\else 0\fi}
\newcommand{\CombinationID}{%
  \ifnum\value{combination}<10 0\fi\arabic{combination}%
}

\newcommand{\MobilityRoute}[3]{%
  \ifnum#1=1
    M3*
  \else
    \ifnum#2=1
      M3*
    \else
      \ifnum#3=1 M2*\else M2/M1*\fi
    \fi
  \fi
}

\newcommand{\FunctionRoute}[3]{%
  \ifnum#1=1
    \ifnum#2=1
      U3*
    \else
      \ifnum#3=1 U3/U2*\else U2*\fi
    \fi
  \else
    \ifnum#2=1
      U2/U1*
    \else
      \ifnum#3=1 U2/U1*\else U1*\fi
    \fi
  \fi
}

\newcommand{\RegistryRow}[7]{%
  #1 & \Door{#2} & \Door{#3} & \Door{#4} & \Door{#5} & \Door{#6} & \Door{#7} &
  \MobilityRoute{#5}{#6}{#3} & \FunctionRoute{#2}{#4}{#7} \\
}

\begin{document}

\begin{titlepage}
  \centering
  \vspace*{1.2cm}
  {\Large\bfseries SKENARIO AKSES DATA DAN KELUARGA METODOLOGIS\\[8pt]}
  {\large Revisi substantif: 16 keluarga desain dan registri 64 kombinasi akses\\[18pt]}
  {\normalsize \JudulSementara\\[28pt]}
  \begin{tabular}{rl}
    Mahasiswa & \NamaMahasiswa \\
    NIM & \NIMMahasiswa \\
    Program studi & Sarjana Perencanaan Wilayah dan Kota \\
    Departemen & Teknik Arsitektur dan Perencanaan \\
    Fakultas & Teknik, Universitas Gadjah Mada \\
    Pembimbing & \NamaPembimbing
  \end{tabular}
  \vfill
  {\large Yogyakarta\\2026}
\end{titlepage}

\pagenumbering{roman}
\tableofcontents
\clearpage
\pagenumbering{arabic}
\hypersetup{pageanchor=true}

\section{Status dan fungsi dokumen}

Dokumen ini adalah rancangan kerja sebelum keluaran permohonan data diperiksa. Ia tidak menyatakan bahwa suatu data telah diterima, lengkap, atau layak dianalisis. Enam jawaban akses tetap diregistrasikan secara biner untuk keperluan administrasi, tetapi desain penelitian tidak lagi ditentukan langsung dari angka 1 atau 0.

Revisi ini memisahkan tiga hal yang sebelumnya tercampur:
\begin{enumerate}
  \item \textbf{akses permohonan}, yaitu apakah enam pintu data menghasilkan berkas yang dapat diperiksa;
  \item \textbf{mutu bukti}, yaitu apakah keluaran tersebut langsung, agregat, sintetis, proksi, atau harus dikeluarkan; dan
  \item \textbf{keluarga metodologis}, yaitu desain analisis serta batas klaim yang dibolehkan oleh kombinasi bukti integrasi dan fungsi lokal.
\end{enumerate}

Seluruh 64 kombinasi biner dipertahankan dalam Lampiran A sebagai registri audit akses. Kombinasi itu bukan 64 desain substantif yang berbeda. Bagian utama hanya mempertahankan keluarga metodologis dan skenario prioritas yang benar-benar mengubah indikator, unit, metode, atau klaim.

\section{Keputusan desain yang dipertahankan}

\begin{itemize}
  \item Kasus utama adalah Jabodetabek/BPTJ, sedangkan Wilayah Statistik Metropolitan (WSM) menjadi sabuk pembanding fungsional.
  \item DKI Jakarta diperlakukan sebagai satu inti untuk hubungan eksternal dengan batas administratif tetap.
  \item Baseline dominan adalah 2024. Data 2023 atau periode lain tidak dilebur seolah-olah berasal dari tahun yang sama; tahun setiap indikator tetap ditampilkan.
  \item Desain bersifat multiresolusi: kecamatan menjadi unit utama jika data kritis mendukung, titik atau raster menjadi lapisan pendukung, dan kabupaten/kota menjadi konteks.
  \item Resolusi analisis mengikuti sumber kritis yang paling kasar. Data agregat tidak diturunkan secara artifisial ke kecamatan, zona, grid, atau titik.
  \item Baseline sumber terbuka dan konfigurasi data terbatas diperlakukan sebagai dua hasil yang setara dan dapat dibandingkan, bukan sebagai versi publik yang dianggap inferior secara otomatis.
  \item Tipologi deskriptif dapat dibangun pada seluruh keluarga yang masih layak. Label ketat \textit{borrowed size} atau \textit{agglomeration shadow} hanya digunakan setelah gerbang bukti relasional dan benchmark fungsi terpenuhi.
  \item Skenario bukti dan skenario parameter WebGL tetap dipisahkan. Perubahan sumber data tidak boleh ditampilkan sebagai perubahan kondisi metropolitan.
\end{itemize}

\section{Pintu akses dan baseline terbuka}

\subsection{Enam pintu permohonan}

\begin{tabularx}{\textwidth}{@{}P{0.07\textwidth}P{0.28\textwidth}Y@{}}
\toprule
Kode & Pintu data & Kontribusi yang diminta \\
\midrule
P & Podes 2024 & Kehadiran dan akses layanan lokal pada desa/kelurahan yang dapat diagregasikan secara sah. \\
C & Survei Komuter Jabodetabek 2023 & OD, karakteristik komuter, dan beban perjalanan pada resolusi diseminasi yang sah. \\
S & Sakernas Agustus 2024 & Pekerjaan menurut lokasi tempat kerja dan sektor luas pada unit dengan reliabilitas yang dapat diterima. \\
D & Kemenhub/DJITM & Matriks atau tabulasi OD Kajian Transportasi Jabodetabek 2023 beserta zona dan metadata pembentukannya. \\
J & JUTPI 2/3 & OD antarzona, TAZ/crosswalk, jaringan, waktu tempuh, serta informasi kalibrasi dan validasi. \\
G & Disnaker/DPMPTSP & Tabulasi pekerjaan atau fungsi ekonomi pada 3--5 pusat sekunder prioritas untuk triangulasi lokal. \\
\bottomrule
\end{tabularx}

Kode G hanya menyederhanakan pencatatan permohonan. Setelah jawaban diterima, Disnaker dan DPMPTSP harus dicatat terpisah karena cakupan, definisi, unit, dan kelengkapannya dapat berbeda. Persetujuan salah satu instansi tidak boleh dianggap sebagai cakupan sistem metropolitan secara keseluruhan.

\subsection{Baseline yang tidak bergantung pada enam permohonan}

Kegagalan enam pintu tidak berarti seluruh bukti bernilai nol. Baseline terbuka tetap dapat memuat WSM 2024, tabel marginal bangkitan--tarikan BPTJ, jaringan dan GTFS yang tersedia, batas administrasi, penduduk, PDRB kabupaten/kota sebagai konteks, fasilitas resmi yang dapat diperiksa, serta proksi OSM/Overture, morfologi terbangun, dan sumber terbuka lain yang sudah diterima dalam rancangan penelitian.

Namun, baseline tersebut tetap tunduk pada audit. Persentase orientasi komuter WSM tidak sama dengan matriks OD antarpusat; bangkitan--tarikan tidak sama dengan pasangan arus; jaringan mengukur aksesibilitas potensial; dan POI tidak sama dengan sensus fungsi atau pekerjaan.

\section{Status mutu bukti}

Setelah berkas diterima, setiap modul diberi status berikut. Status ditentukan dari isi berkas, bukan dari nama produk atau surat persetujuan.

\begin{tabularx}{\textwidth}{@{}P{0.08\textwidth}P{0.18\textwidth}Y@{}}
\toprule
Status & Nama & Aturan penggunaan \\
\midrule
D & Langsung & Konstruk teramati pada unit dan periode sasaran dengan definisi serta kualitas yang memadai. \\
A & Agregat & Konstruk teramati, tetapi wilayah, waktu, kategori, atau cakupannya lebih kasar daripada unit utama. \\
S & Sintetis & Konstruk diestimasi melalui model dari data lain; asumsi, kendala, kalibrasi, dan ketidakpastian wajib ditampilkan. \\
P & Proksi & Sumber mengukur gejala terkait, bukan konstruk yang hendak diklaim secara langsung. \\
$\varnothing$ & Dikeluarkan & Data tidak tersedia, tidak dapat dibandingkan, tidak reliabel, atau tidak mempunyai pengganti yang dapat dipertanggungjawabkan. \\
\bottomrule
\end{tabularx}

Contoh konsekuensi: keluaran JUTPI yang merupakan hasil model dapat berstatus S meskipun pintu J bernilai 1; Sakernas yang hanya sah pada kabupaten/kota dapat berstatus A untuk penelitian berkecamatan; dan data Pemda yang hanya mencakup pusat terpilih tidak otomatis berstatus D untuk seluruh sistem.

\section{Pertanyaan dan tujuan penelitian}

Pertanyaan substantif dipertahankan stabil. Sumber data menentukan kedalaman jawaban dan kekuatan klaim, bukan mengganti konstruk penelitian dengan nama dataset.

\subsection{Rumusan masalah}

\begin{enumerate}
  \item Bagaimana derajat, arah, dan beban integrasi fungsional pusat-pusat sekunder dengan Jakarta dan, jika data mendukung, dengan pusat lain berbeda secara spasial?
  \item Bagaimana fungsi dan kinerja lokal setiap pusat dibandingkan dengan tingkat yang diharapkan berdasarkan ukuran serta karakteristik lokalnya?
  \item Bagaimana kombinasi integrasi, surplus atau defisit fungsi, manfaat, dan beban ketergantungan membentuk tipologi serta gradien kondisi relasional metropolitan pada status quo?
  \item Seberapa stabil tipologi dan jejak spasial tersebut terhadap perubahan indikator, unit, bobot, ambang, dan status mutu bukti, dengan batas administratif tetap?
\end{enumerate}

Jika suatu keluarga tidak memiliki bukti arus, pertanyaan pertama hanya dijawab sebagai aksesibilitas potensial. Jika benchmark fungsi tidak dapat dibangun, pertanyaan kedua hanya dijawab sebagai profil fungsi atau proksi. Dalam dua keadaan tersebut, pertanyaan ketiga tidak menghasilkan diagnosis ketat \textit{borrowed size--shadow}.

\subsection{Tujuan penelitian}

\begin{enumerate}
  \item Mengukur dan memetakan integrasi fungsional, orientasi perjalanan, atau aksesibilitas potensial sesuai status bukti yang tersedia.
  \item Menilai kematangan dan fungsi relatif pusat sekunder melalui benchmark yang sesuai dengan unit dan konstruk yang benar-benar teramati.
  \item Menyusun tipologi dan gradien status quo dengan membedakan hasil deskriptif dari diagnosis konseptual yang memenuhi gerbang bukti.
  \item Menguji ketahanan hasil terhadap perubahan indikator, unit, bobot, ambang, serta konfigurasi bukti melalui analisis sensitivitas dan penjelajah WebGL.
\end{enumerate}

\section{Dua sumbu keluarga metodologis}

\subsection{Sumbu M: bukti mobilitas dan hubungan metropolitan}

\begin{tabularx}{\textwidth}{@{}P{0.08\textwidth}P{0.24\textwidth}Y@{}}
\toprule
Kode & Tingkat bukti & Metode dan batas \\
\midrule
M3 & OD relasional rinci & OD antarkecamatan/TAZ atau antarpusat tersedia. Ukur orientasi ke inti, hubungan antarpusat, inflow--outflow, dan self-containment hanya jika definisi matriks memungkinkan. Bedakan arus teramati dan termodelkan. \\
M2 & Relasi terarah agregat & Persentase orientasi ke inti, OD kabupaten/kota, atau tabulasi komuter tersedia. Analisis sah pada resolusi sumber dan dapat mengukur beban komuter, tetapi tidak mengarang hubungan antarkecamatan atau antarpusat. \\
M1 & Sintetis atau aksesibilitas & OD sintetis, marginal bangkitan--tarikan, jaringan, atau waktu tempuh digunakan. Keluaran disebut potensi interaksi atau aksesibilitas, bukan integrasi aktual. \\
M0 & Tidak ada bukti layak & Tidak tersedia relasi, model, atau proksi jaringan yang dapat dipertanggungjawabkan. Modul integrasi dikeluarkan dan tesis inti tidak dapat didiagnosis. \\
\bottomrule
\end{tabularx}

Jika D, J, dan C tersedia bersamaan, ketiganya tidak disusun sebagai hierarki yang membuang sumber lain. D/J digunakan untuk struktur OD pada unitnya; C digunakan untuk karakteristik individu dan beban perjalanan; WSM dan sumber lain menjadi pemeriksaan silang sesuai cakupan dan tahun.

\subsection{Sumbu U: bukti kematangan dan fungsi lokal}

\begin{tabularx}{\textwidth}{@{}P{0.08\textwidth}P{0.24\textwidth}Y@{}}
\toprule
Kode & Tingkat bukti & Metode dan batas \\
\midrule
U3 & Layanan dan pekerjaan kompatibel & Podes atau bukti layanan setara tersedia bersama pekerjaan menurut lokasi kerja pada unit yang kompatibel. Benchmark dapat mencakup kematangan layanan dan ekonomi, dengan reliabilitas serta cakupan dilaporkan. \\
U2 & Satu dimensi langsung & Layanan lokal teramati tetapi pekerjaan hanya agregat/tidak tersedia, atau pekerjaan teramati tetapi layanan tidak setara Podes. Benchmark dibatasi pada dimensi yang terukur; kematangan ekonomi penuh tidak diklaim. \\
U1 & Fungsi berbasis proksi & POI, direktori, morfologi terbangun, kawasan industri, atau konteks agregat digunakan. Keluaran disebut kehadiran/proksi fungsi dan diaudit terhadap cakupan serta duplikasi. \\
U0 & Tidak ada bukti layak & Tidak tersedia indikator fungsi yang dapat dibandingkan. Modul fungsi relatif dikeluarkan dan diagnosis \textit{borrowed size--shadow} tidak dapat dibangun. \\
\bottomrule
\end{tabularx}

Jika S dan G tersedia bersamaan, data Pemda dipakai untuk validasi lokal, bukan dihilangkan oleh prioritas Sakernas. Hasil G hanya digeneralisasikan kepada pusat dan cakupan pelaporan yang benar-benar dilayani.

\section{Enam belas keluarga desain}

\begin{longtable}{@{}P{0.10\textwidth}P{0.13\textwidth}P{0.25\textwidth}P{0.44\textwidth}@{}}
\toprule
Keluarga & Tingkat & Keluaran yang sah & Batas klaim utama \\
\midrule
\endfirsthead
\toprule
Keluarga & Tingkat & Keluaran yang sah & Batas klaim utama \\
\midrule
\endhead
M3--U3 & A & Diagnosis relasional dan benchmark fungsi lengkap & Kandidat penggunaan label ketat, hanya setelah unit, periode, benchmark, dan provenance lolos audit. \\
M3--U2 & B & Arus rinci dan surplus/defisit layanan atau ekonomi parsial & Pola dapat disebut konsisten dengan tesis, tetapi kematangan fungsi tidak boleh digeneralisasikan ke dimensi yang hilang. \\
M3--U1 & C & Profil arus rinci disandingkan dengan proksi fungsi & Tidak ada diagnosis ketat; hasil utama adalah ketimpangan arus dan profil fungsi proksi. \\
M3--U0 & C & Atlas arus dan beban hubungan metropolitan & Tidak dapat membedakan penguatan fungsi dari ketergantungan lokal. \\
M2--U3 & A* & Diagnosis terarah dengan fungsi lokal kuat & Label ketat hanya jika relasi terarah dan fungsi berada pada unit/catchment yang sebanding; jika tidak, turun ke B. \\
M2--U2 & B & Desain minimum integrasi--kematangan layanan & Mendukung pola yang konsisten dengan kondisi produktif, campuran, atau ketergantungan; bukan diagnosis ekonomi penuh. \\
M2--U1 & C & Orientasi komuter dan fungsi proksi & Tidak boleh mengklaim surplus produktivitas, self-containment pusat, atau shadow antarpusat. \\
M2--U0 & C & Profil orientasi dan beban komuter & Tidak dapat menjawab kematangan fungsi lokal. \\
M1--U3 & C & Fungsi lokal kuat dengan aksesibilitas/potensi interaksi & Tidak boleh menyebut integrasi aktual atau mengatribusikan surplus/defisit fungsi kepada hubungan dengan Jakarta. \\
M1--U2 & C & Benchmark fungsi parsial dan aksesibilitas potensial & Hasil berupa profil aksesibilitas--fungsi, bukan diagnosis relasional. \\
M1--U1 & D & Baseline proksi multidimensi & Berguna untuk pemetaan awal dan stress test; label konseptual ketat dilarang. \\
M1--U0 & D & Aksesibilitas potensial saja & Tesis fungsi lokal tidak terjawab. \\
M0--U3 & C & Atlas fungsi dan benchmark lokal & Tidak dapat mengaitkan kondisi fungsi dengan integrasi Jakarta. \\
M0--U2 & C & Profil fungsi lokal parsial & Tidak dapat menguji tesis relasional. \\
M0--U1 & D & Inventaris fungsi proksi & Hanya eksplorasi deskriptif. \\
M0--U0 & E & Tidak ada keluaran inti & Penelitian harus menunggu data, mengubah unit, atau merancang ulang pertanyaan. \\
\bottomrule
\end{longtable}

Tingkat A adalah kandidat diagnosis penuh; B adalah versi minimum yang masih menjawab tesis secara terbatas; C hanya menjawab satu sisi atau hubungan parsial; D adalah baseline proksi; dan E tidak layak untuk pertanyaan inti. Tanda bintang menunjukkan bahwa kelayakan sangat bergantung pada kompatibilitas unit.

\section{Metode bersama dan gerbang klaim}

\subsection{Audit dan harmonisasi}

Setiap berkas dicatat menurut sumber, versi, tahun observasi, tanggal akses, unit, cakupan, definisi, metode pembentukan, missingness, reliabilitas, lisensi, dan pembatasan publikasi. Baseline 2024 menjadi jangkar; indikator 2023 atau periode lain dipertahankan sebagai pengamatan berlabel, bukan dirata-ratakan tanpa alasan.

Harmonisasi spasial mengikuti urutan berikut:
\begin{enumerate}
  \item gunakan unit asli untuk perhitungan awal;
  \item gunakan crosswalk resmi jika tersedia;
  \item agregasikan unit yang lebih rinci ke unit pembanding yang lebih kasar jika operasi tersebut sah;
  \item jika unit tidak dapat disetarakan, tampilkan analisis paralel multiresolusi;
  \item jangan mendisagregasikan nilai kabupaten/kota atau zona kasar menjadi variasi faktual kecamatan/grid.
\end{enumerate}

\subsection{Pengukuran integrasi}

M3 dapat menghitung orientasi ke inti, keterkaitan antarpusat, inflow--outflow, dan self-containment jika definisi OD mendukung. M2 membatasi analisis pada orientasi dan beban pada unit sumber. M1 memodelkan potensi aksesibilitas atau interaksi dengan asumsi yang terlihat. Tidak satu pun keluaran model disebut arus teramati.

\subsection{Pengukuran fungsi relatif}

Indikator layanan, pekerjaan, fungsi orde tinggi, dan kegiatan ekonomi tidak dilebur sebelum provenance serta polaritasnya diperiksa. Fungsi yang teramati dibandingkan dengan tingkat yang diharapkan dari ukuran dan karakteristik lokal melalui benchmark transparan atau model fungsi--ukuran. Residual positif atau negatif hanya berarti surplus atau defisit relatif pada konstruk yang dimodelkan. Residual fasilitas tidak disebut surplus produktivitas, dan data potong lintang tidak disebut \textit{functional upgrading} temporal.

\subsection{Tipologi dan sensitivitas}

Manfaat akses, kematangan fungsi, dan beban ketergantungan dipertahankan sebagai dimensi terpisah. Kategori campuran tetap tersedia. Hasil dilaporkan sebagai gradien sebelum diringkas menjadi kelas. Uji ketahanan sekurang-kurangnya memeriksa perubahan unit, korelasi dan duplikasi indikator, \textit{leave-one-indicator-out}, bobot, ambang, serta perubahan status bukti.

\subsection{Empat gerbang penggunaan label ketat}

Label \textit{borrowed size} atau \textit{agglomeration shadow} hanya dapat dipromosikan jika seluruh syarat berikut terpenuhi:
\begin{enumerate}
  \item tersedia bukti relasional terhadap Jakarta, bukan hanya kedekatan atau waktu tempuh;
  \item tersedia benchmark fungsi relatif yang mengukur konstruk substantif pada unit yang memadai;
  \item bukti relasional dan fungsi dapat dihubungkan pada unit, catchment, serta periode yang dapat dipertanggungjawabkan; dan
  \item klasifikasi tidak runtuh pada uji sensitivitas utama dan tidak bergantung pada satu proksi bermasalah.
\end{enumerate}

Jika satu gerbang gagal, hasil tetap dilaporkan sebagai profil atau pola yang konsisten dengan konstruk tertentu, dengan alasan penurunan klaim dinyatakan.

\section{Skenario prioritas}

Tabel berikut menggantikan uraian berulang untuk 64 kombinasi. Kode keluarga masih bersifat awal karena status akhir baru dapat ditentukan setelah berkas diperiksa.

\begin{longtable}{@{}P{0.09\textwidth}P{0.11\textwidth}P{0.11\textwidth}P{0.21\textwidth}P{0.31\textwidth}@{}}
\toprule
ID & Vektor & Rute awal & Alasan dipertahankan & Konsekuensi substantif \\
\midrule
\endfirsthead
\toprule
ID & Vektor & Rute awal & Alasan dipertahankan & Konsekuensi substantif \\
\midrule
\endhead
EV-FULL & 111111 & M3*/U3* & Semua pintu menghasilkan keluaran & D dan J dibandingkan, C mengukur karakteristik/beban, sedangkan S dan G saling memeriksa. Tidak ada sumber yang diabaikan hanya karena sumber lain lebih rinci. \\
EV-P0 & 011111 & M3*/U2--U1* & Kegagalan tunggal Podes & Fungsi layanan memakai proksi; pekerjaan tidak mengembalikan informasi layanan yang hilang. Diagnosis ekonomi hanya berlaku pada unit pekerjaan yang sah. \\
EV-C0 & 101111 & M3*/U3* & Kegagalan tunggal Komuter & Struktur OD masih mungkin kuat, tetapi karakteristik individu, biaya, moda, dan beban survei tidak dapat diklaim dari D/J tanpa variabel setara. \\
EV-S0 & 110111 & M3*/U3--U2* & Kegagalan tunggal Sakernas & G menjadi validasi pusat terpilih, bukan pengganti sistem metropolitan. Jika cakupan G sempit, fungsi ekonomi turun ke U2. \\
EV-D0 & 111011 & M3*/U3* & Kegagalan tunggal DJITM & M3 hanya aktif jika J memuat OD dan zona yang dapat dipakai; jika hanya metadata/geometri, turun ke M2 atau M1. \\
EV-J0 & 111101 & M3*/U3* & Kegagalan tunggal JUTPI & M3 hanya aktif jika keluaran D benar-benar matriks relasional, bukan sekadar marginal bangkitan--tarikan. \\
EV-G0 & 111110 & M3*/U3--U2* & Kegagalan tunggal Pemda & S dapat menopang pekerjaan hanya pada unit dan tingkat reliabilitas yang sah; validasi pusat terpilih hilang. \\
EV-BPS & 111000 & M2*/U3--U2* & Hanya tiga pintu BPS terbuka & C tetap agregat pada resolusi sumber. P dan S tidak boleh dipadankan pada kecamatan jika S hanya sah pada kabupaten/kota. \\
EV-MIN & 110000 & M2*/U2* & Paket minimum Podes--Komuter & Menjawab orientasi/beban dan kematangan layanan; hanya mendukung pola konsisten, bukan kematangan ekonomi penuh. \\
EV-OD & 101100 & M3*/U3--U2* & Podes, Sakernas, dan DJITM & Kandidat desain kuat jika OD dan pekerjaan kompatibel. Jika Sakernas agregat, analisis menjadi multiresolusi dan turun ke U2. \\
EV-JG & 100011 & M3*/U3--U2* & Alternatif JUTPI dan Pemda & Diagnosis hanya untuk pusat yang tercakup G serta zona JUTPI yang dapat dihubungkan secara sah. \\
EV-FLOW & 010110 & M3*/U1* & Sumber mobilitas terbuka, fungsi resmi tertutup & Menilai struktur arus dan beban dengan fungsi proksi. Label ketat dilarang karena benchmark fungsi substantif tidak tersedia. \\
EV-FUNC & 101001 & M2 atau M1*; U3 atau U2* & Sumber fungsi terbuka, OD permohonan tertutup & WSM/baseline terbuka menentukan M2 atau M1. Surplus fungsi tidak boleh diatribusikan kepada integrasi Jakarta tanpa bukti relasional memadai. \\
EV-PODES & 100000 & M2 atau M1*; U2* & Hanya Podes yang diperoleh & Menghasilkan benchmark layanan dan profil aksesibilitas/orientasi terbatas, bukan fungsi ekonomi atau perubahan temporal. \\
EV-OPEN & 000000 & M2 atau M1*; U1* & Seluruh permohonan gagal & Baseline terbuka tetap berjalan. Keluaran adalah profil orientasi/aksesibilitas dan fungsi proksi; tidak ada diagnosis ketat. \\
EV-AKT & diisi kemudian & ditetapkan setelah audit & Kombinasi yang benar-benar terjadi & Manifest aktual menjadi skenario utama bersama baseline terbuka. Keluarga dipilih dari isi dan mutu keluaran, bukan jumlah pintu terbuka. \\
\bottomrule
\end{longtable}

Skenario lain ditambahkan hanya jika benar-benar terjadi atau menghasilkan perubahan metode/klaim yang tidak terwakili oleh tabel. Dua kombinasi dengan keluarga, indikator, unit, dan batas klaim yang sama tidak perlu ditulis sebagai dua desain penuh.

\section{Manifest skenario aktual}

Ketika jawaban instansi diterima, satu manifest diisi untuk setiap skenario yang dipertahankan.

\begin{longtable}{@{}P{0.09\textwidth}P{0.18\textwidth}P{0.12\textwidth}P{0.18\textwidth}P{0.29\textwidth}@{}}
\toprule
Modul & Konstruk & Status & Sumber, tahun, unit & Indikator dan batas \\
\midrule
F & Arus dan orientasi & D/A/S/P/$\varnothing$ & diisi setelah audit & OD, orientasi inti, atau aksesibilitas; jangan campur status tanpa label. \\
Jb & Pekerjaan tempat kerja & D/A/S/P/$\varnothing$ & diisi setelah audit & Jumlah/sektor pekerjaan dan reliabilitas unit. \\
L & Layanan lokal & D/A/S/P/$\varnothing$ & diisi setelah audit & Kehadiran/akses layanan; bukan produktivitas. \\
E & Usaha serta HQ--cabang & D/A/S/P/$\varnothing$ & diisi setelah audit & Kehadiran/keragaman usaha; status kantor hanya jika teramati. \\
R & Output, upah, produktivitas & D/A/S/P/$\varnothing$ & diisi setelah audit & Konteks agregat tidak diturunkan ke kecamatan. \\
N & Jaringan/waktu tempuh & D/A/S/P/$\varnothing$ & diisi setelah audit & Aksesibilitas potensial; bukan arus penumpang. \\
H & Layanan orde tinggi & D/A/S/P/$\varnothing$ & diisi setelah audit & Kehadiran, kelas, kapasitas, dan status verifikasi. \\
I & Struktur industri & D/A/S/P/$\varnothing$ & diisi setelah audit & Poligon tidak membuktikan tenant, pekerjaan, atau upgrading. \\
K & Lahan/perumahan & D/A/S/P/$\varnothing$ & diisi setelah audit & Data pengguna belum dianggap terverifikasi sebelum berkas diperiksa. \\
T & Perubahan temporal & D/A/S/P/$\varnothing$ & diisi setelah audit & Tanpa seri harmonis, keluaran hanya status quo. \\
Gv & Batas/tanggung jawab & D/A/S/P/$\varnothing$ & diisi setelah audit & Overlay administratif tetap; bukan hasil perluasan Jakarta. \\
\bottomrule
\end{longtable}

Manifest juga mencatat indikator yang boleh dihitung, klaim yang dilarang, uji validasi, lapisan WebGL, aturan agregasi, versi berkas, serta reproduksi pengolahan. Jika satu berkas berubah versi atau status, skenario diberi versi baru tanpa menimpa provenance lama.

\section{Aturan WebGL}

Versi publik hanya memuat data terbuka, agregat, atau hasil turunan yang diizinkan. Versi penelitian lokal dapat memakai data terbatas sesuai perjanjian. Antarmuka selalu menampilkan ID skenario bukti, status setiap modul, unit, tahun, dan penanda proksi/sintetis.

Slider bobot, ambang, atau nilai indikator tidak boleh:
\begin{itemize}
  \item mengubah data agregat menjadi variasi faktual pada unit yang lebih rinci;
  \item memperlakukan proksi sebagai pengamatan langsung;
  \item menyembunyikan perubahan status bukti melalui normalisasi ulang; atau
  \item menggambarkan perubahan kelas sebagai ramalan kebijakan, perubahan yurisdiksi, atau hubungan kausal.
\end{itemize}

\clearpage
\appendix
\section{Registri 64 kombinasi akses}

Lampiran ini hanya menunjukkan ruang kemungkinan jawaban enam permohonan. Tanda bintang berarti rute keluarga masih sementara. Keluarga final ditentukan setelah isi, unit, periode, reliabilitas, dan provenance keluaran diperiksa. Kolom M memakai urutan awal D/J, C, lalu baseline terbuka; kolom U memakai kombinasi P, S, G, lalu proksi terbuka.

{\small
\setlength{\tabcolsep}{4pt}
\begin{longtable}{@{}P{0.07\textwidth}ccccccP{0.12\textwidth}P{0.14\textwidth}@{}}
\toprule
No. & P & C & S & D & J & G & M awal & U awal \\
\midrule
\endfirsthead
\toprule
No. & P & C & S & D & J & G & M awal & U awal \\
\midrule
\endhead
\RegistryRow{01}{1}{1}{1}{1}{1}{1}
\RegistryRow{02}{1}{1}{1}{1}{1}{0}
\RegistryRow{03}{1}{1}{1}{1}{0}{1}
\RegistryRow{04}{1}{1}{1}{1}{0}{0}
\RegistryRow{05}{1}{1}{1}{0}{1}{1}
\RegistryRow{06}{1}{1}{1}{0}{1}{0}
\RegistryRow{07}{1}{1}{1}{0}{0}{1}
\RegistryRow{08}{1}{1}{1}{0}{0}{0}
\RegistryRow{09}{1}{1}{0}{1}{1}{1}
\RegistryRow{10}{1}{1}{0}{1}{1}{0}
\RegistryRow{11}{1}{1}{0}{1}{0}{1}
\RegistryRow{12}{1}{1}{0}{1}{0}{0}
\RegistryRow{13}{1}{1}{0}{0}{1}{1}
\RegistryRow{14}{1}{1}{0}{0}{1}{0}
\RegistryRow{15}{1}{1}{0}{0}{0}{1}
\RegistryRow{16}{1}{1}{0}{0}{0}{0}
\RegistryRow{17}{1}{0}{1}{1}{1}{1}
\RegistryRow{18}{1}{0}{1}{1}{1}{0}
\RegistryRow{19}{1}{0}{1}{1}{0}{1}
\RegistryRow{20}{1}{0}{1}{1}{0}{0}
\RegistryRow{21}{1}{0}{1}{0}{1}{1}
\RegistryRow{22}{1}{0}{1}{0}{1}{0}
\RegistryRow{23}{1}{0}{1}{0}{0}{1}
\RegistryRow{24}{1}{0}{1}{0}{0}{0}
\RegistryRow{25}{1}{0}{0}{1}{1}{1}
\RegistryRow{26}{1}{0}{0}{1}{1}{0}
\RegistryRow{27}{1}{0}{0}{1}{0}{1}
\RegistryRow{28}{1}{0}{0}{1}{0}{0}
\RegistryRow{29}{1}{0}{0}{0}{1}{1}
\RegistryRow{30}{1}{0}{0}{0}{1}{0}
\RegistryRow{31}{1}{0}{0}{0}{0}{1}
\RegistryRow{32}{1}{0}{0}{0}{0}{0}
\RegistryRow{33}{0}{1}{1}{1}{1}{1}
\RegistryRow{34}{0}{1}{1}{1}{1}{0}
\RegistryRow{35}{0}{1}{1}{1}{0}{1}
\RegistryRow{36}{0}{1}{1}{1}{0}{0}
\RegistryRow{37}{0}{1}{1}{0}{1}{1}
\RegistryRow{38}{0}{1}{1}{0}{1}{0}
\RegistryRow{39}{0}{1}{1}{0}{0}{1}
\RegistryRow{40}{0}{1}{1}{0}{0}{0}
\RegistryRow{41}{0}{1}{0}{1}{1}{1}
\RegistryRow{42}{0}{1}{0}{1}{1}{0}
\RegistryRow{43}{0}{1}{0}{1}{0}{1}
\RegistryRow{44}{0}{1}{0}{1}{0}{0}
\RegistryRow{45}{0}{1}{0}{0}{1}{1}
\RegistryRow{46}{0}{1}{0}{0}{1}{0}
\RegistryRow{47}{0}{1}{0}{0}{0}{1}
\RegistryRow{48}{0}{1}{0}{0}{0}{0}
\RegistryRow{49}{0}{0}{1}{1}{1}{1}
\RegistryRow{50}{0}{0}{1}{1}{1}{0}
\RegistryRow{51}{0}{0}{1}{1}{0}{1}
\RegistryRow{52}{0}{0}{1}{1}{0}{0}
\RegistryRow{53}{0}{0}{1}{0}{1}{1}
\RegistryRow{54}{0}{0}{1}{0}{1}{0}
\RegistryRow{55}{0}{0}{1}{0}{0}{1}
\RegistryRow{56}{0}{0}{1}{0}{0}{0}
\RegistryRow{57}{0}{0}{0}{1}{1}{1}
\RegistryRow{58}{0}{0}{0}{1}{1}{0}
\RegistryRow{59}{0}{0}{0}{1}{0}{1}
\RegistryRow{60}{0}{0}{0}{1}{0}{0}
\RegistryRow{61}{0}{0}{0}{0}{1}{1}
\RegistryRow{62}{0}{0}{0}{0}{1}{0}
\RegistryRow{63}{0}{0}{0}{0}{0}{1}
\RegistryRow{64}{0}{0}{0}{0}{0}{0}
\bottomrule
\end{longtable}
}

\section{Checklist aktivasi skenario}

\begin{enumerate}
  \item Catat jawaban enam pintu tanpa menganggap persetujuan sebagai mutu data.
  \item Inventarisasi berkas aktual dan pisahkan jawaban Disnaker dari DPMPTSP.
  \item Beri status D/A/S/P/$\varnothing$ kepada setiap modul berdasarkan isi berkas.
  \item Tetapkan M0--M3 dan U0--U3, termasuk aturan unit bersama atau analisis paralel multiresolusi.
  \item Pilih EV-AKT, baseline terbuka, kegagalan tunggal yang relevan, dan stress test yang mengubah klaim.
  \item Nyatakan indikator yang dihapus, diganti, atau diturunkan resolusinya; jangan mengimputasi data hilang sebagai nol.
  \item Terapkan empat gerbang klaim sebelum menggunakan label konseptual ketat.
  \item Simpan konfigurasi bukti secara terpisah dari variasi parameter WebGL.
\end{enumerate}

\end{document}
